%%%%% TEXTOS SOBRE AGRICULTURA SUSTENTÁVEL



Os diversos conceitos relacionam a propriedade e seu gerenciamento, o uso de mão-de-obra e tamanho físico ou econômico da mesma. Existe um consenso de que a propriedade rural é operada, ou gerida, por um membro de um agregado familiar e que o trabalho é realizado pelo proprietário e sua família \cite{Schneider2003}. Em 2014, o Comitê Internacional de Direção para o Ano Internacional da Agricultura Familiar, a definiu para organizar a agricultura, silvicultura, pesca, pastoral e produção aquícola gerido e operado por uma família e predominantemente dependente do trabalho familiar, incluindo mulheres e homens. Nestes sistemas, a família e a fazenda estão ligadas, co-evoluem: economia, meio ambiente, funções sociais e culturais \cite{FAO2014}.

Em 2018, o Instituto Brasileiro de Geografia e Estatística (IBGE), em parceria com o Ministério do Desenvolvimento Agrário (MDA), publicou a parcial do Censo Agropecuário 2017 \cite{IBGE2018CensoParanaense} relatando que no Brasil atualmente, 23\% das propriedades rurais são de agricultores familiares, ratificando a importância da Agricultura Familiar no Brasil. A agricultura familiar é fundamental para a provisão de alimentos base para a constituição de dietas saudáveis, tais como legumes, frutas, vegetais e de origem animal. Apesar de os 4.367.902 de estabelecimentos familiares representarem 84,4\% do total, estas propriedades ocupam apenas 24,3\% da área ocupada pelos empreendimentos agropecuários brasileiros. 

Confirmando os estudos da FAO, a agricultura familiar é fundamental para a  segurança alimentar do Brasil, responsável por 87,0\% da produção nacional de mandioca; 70,0\% da produção de feijão, 21,0\% da produção do trigo e 16,0\% da soja, principais produtos da pauta de exportação brasileira. Essa modalidade de propriedade também é responsável por 59,0\% do plantel de suínos, 50,0\% do plantel de aves, 30,0\% dos bovinos, \cite{IBGE2018CensoParanaense}.

Definir o desenvolvimento rural sustentável requer um esforço observacional e prático, pois, este ambiente vem sofrendo profundas transformações em suas demandas e necessidades. O desenvolvimento que antes se apresentava majoritariamente como produção de subsistência, hoje dá lugar a um complexo sistema agroindustrial \cite{Bastos2018DeterminantesBrasileiro} e social. É importante neste sentindo compreender que definir o desenvolvimento rural com apenas um conceito seria uma proposição simplista do contexto de crescimento no meio rural. Partindo da definição de consequência de ações governamentais definidas por \citeonline{Navarro2001DesenvolvimentoFuturo} como “atitudes práticas”, este autor descreve que:

\begin{citacao}
[...] Desenvolvimento rural, portanto, pode ser analisado a “posteriori”, neste caso se referindo às análises sobre programas já realizados pelo Estado (em seus diferentes níveis) visando a alterar facetas do mundo rural a partir de objetivos previamente definidos. Porem pode se referir também à elaboração de uma ação prática.
\end{citacao}


Neste campo econômico, a agricultura sustentável, especialmente a agricultura familiar, é um componente indispensável do ponto de vista econômico, social e ambiental. Mesmo no mundo de hoje, que enfrenta desafios como superpopulação, mudanças climáticas, dificuldades na produção agrícola em grande escala e competição de recursos, esta atividade desempenha um papel único \cite{Riegler2018,Rockstrom2017}.

Ao atendar aos requisitos do desenvolvimento sustentável, a agricultura familiar busca eliminar a fome no meio rural e ajuda os agricultores a gerarem renda, oportunidades de trabalho e a possibilidade de competir com os grandes produtores a partir do cooperativismo e o associativismo, a fim de fornecer alimentos mais saudáveis por culturas mais diversificadas, em que os produtores rurais cultivam várias safras na mesma terra \cite{Santana2014}.

Segundo a \textit{Food and Agriculture Organization of the United Nations} FAO \cite{FAO2017a} a agricultura para consumo inclui todas as atividades agrícolas e familiares e está ligada a várias áreas do desenvolvimento rural, seja a produção de hortaliças e alimentos ou de carne para consumo alimentar. A agricultura familiar é uma forma de organizar a produção agrícola, florestal, pesqueira, pastoril e aquícola gerida essencialmente por familiares.



%%%%% TEXTO SOBRE SCOPE REVIEW  E EMPREENDEDORISMO  


Métodos e inovações para construir o empreendedorismo permitem que todos os setores da economia e da sociedade superem os desafios diários \cite{Cuellar-Galvez2018}. Portanto, a agricultura sustentável precisa usar o empreendedorismo e a inovação empresarial como estratégias para desenvolver sua propriedade ou produção, utilizando todos os recursos disponíveis para melhorar ou até mesmo criar produtos e serviços \cite{Julien2017}.

Atividades de tomada de decisão usando novos métodos em áreas rurais geralmente envolvem extensa pesquisa em uma abundante literatura. Essas atividades são trabalhosas, demoradas e difíceis de serem interpretadas pelos pesquisadores \cite{Xiao2019a}. Nessa perspectiva, parece vantajoso desenvolver mecanismos seguros, automatizados, organizados e controláveis na coleta de evidências e considerar a validade e relevância das afirmações encontradas \cite{Idri2018}, especialmente como a seleção de evidências, trata de novos métodos de produção e produção agrícola. Esse mecanismo deve facilitar aos usuários a otimização do tempo e do trabalho, necessários para manter a confiabilidade dos resultados \cite{Rockstrom2017}.

A  Scope—Systematic Review (SSR) em documentos acadêmicos pode atingir esses objetivos de maneira repetível e controlável, porque tais métodos de revisão mais estruturados são projetados para coletar evidências que atendam aos critérios de elegibilidade pré-especificados para responder a questões de pesquisa específicas \cite{Chandler2021}; especulam sobre as tecnologias mais recentes em tópicos de interesse \cite{Keele2007,Medeiros2015}.

Para identificar, avaliar e interpretar os resultados relevantes das revisões de escopo em trabalhos acadêmicos publicados, conforme os princípios das regras empíricas relacionadas aos levantamentos bibliométricos na literatura, atualmente são amplamente adotados e aceitos os procedimentos de classificação, descrição e qualificação de documentos que adotam as leis: (i) Lei da Produtividade do Pesquisador ou Lotka; (ii) a Lei da Dispersão do Conhecimento Científico ou Lei de Bradford e (iii) a Lei da Distribuição e Frequência das Palavras no Texto ou Lei de Zipf, elaborada em 1926, 1934 e 1949. Tais leis foram projetadas para mapear e compreender o domínio do conhecimento, principalmente considerando os dados gerados pela operação de publicação.

O método de análise do conteúdo trata-se de um conjunto de técnicas aplicadas para a análise de dados qualitativos, que se utilizam da ocorrência de palavras e similitudes dos termos concentrados em corpus de modo a determinar a importância destas, no contexto do assunto abordado \cite{Cassettari2015}. Métodos de Revisão de escopo aplicados às pesquisas sobre empreendedorismo na produção agrícola extensiva podem contribuir para a expansão do campo do conhecimento científico, levantamento de metas, ideias e o momento atual sobre o tema, por conteúdos acadêmicos.

Considerando que o tema empreendedorismo e agricultura sustentável ainda são pouco discutidos em conjunto e que se mostram de suma importância, o presente estudo teve por objetivo demonstrar quais os últimos avanços do empreendedorismo agrícola para o meio agrícola com atividade extensiva por uma revisão de escopo (\textit{Scope—Systematic Review}) da literatura internacional.

%Nessa abordagem sobre ciência e conhecimento científico, encontram-se, nas análises de redes bibliométricas, discussões no sentido da reflexão sobre um dado campo de conhecimento, podendo-se determinar temas de pesquisas, como as publicações estão estruturadas (Klavans & Boyack, 2006) e compreender os alinhamentos entre as redes de pesquisadores (Van Eck & Waltman, 2009). Nesta investigação, considera-se a intenção empreendedora e a sustentabilidade devido, sobretudo, à inexistência de pesquisas dessa natureza, que buscam a relação entre esses temas, e isto pode contribuir diretamente para estruturar mais o campo do conhecimento científico referente ao empreendedorismo, ao identificar e relacionar investigações que envolvem esses fenômenos.
% 
% As publicações científicas evidenciam tendências e influências nas mais diversas áreas do conhecimento, uma vez que são agentes de mudança na ciência e, consequentemente, na compreensão científica. Diante da necessidade de pesquisar como os campos dos conhecimentos sobre intenção empreendedora e sustentabilidade estão se estruturando no âmbito acadêmico, principalmente, considerando as análises de redes bibliométricas, assim como a análise de conteúdo (Knutas et al., 2015), questiona-se: como se estruturam as pesquisas científicas que abordam, conjuntamente, intenção empreendedora e sustentabilidade? Nesse sentido, delimitou-se como objetivo da pesquisa, portanto, investigar a produção científica internacional sobre intenção empreendedora e sustentabilidade. 

% Todavia, por meio de buscas sistêmicas em grandes portais de pesquisas científicas, como Spell, Scielo, Scopus e Web of Science, não foram constatados estudos sobre esses temas concomitantemente com caráter bibliométrico, adiante revistos, razão esta porque se entende que esta pesquisa pode suplantar uma lacuna da literatura científica e, a partir disso, contribuir diretamente para o desenvolvimento dos temas investigados.


%Nessa abordagem sobre ciência e conhecimento científico, encontram-se, nas análises de redes bibliométricas, discussões no sentido da reflexão sobre um dado campo de conhecimento, podendo-se determinar temas de pesquisas, como as publicações estão estruturadas (Klavans & Boyack, 2006) e compreender os alinhamentos entre as redes de pesquisadores (Van Eck & Waltman, 2009). Nesta investigação, considera-se a intenção empreendedora e a sustentabilidade devido, sobretudo, à inexistência de pesquisas dessa natureza, que buscam a relação entre esses temas, e isto pode contribuir diretamente para estruturar mais o campo do conhecimento científico referente ao empreendedorismo, ao identificar e relacionar investigações que envolvem esses fenômenos.
% 
% As publicações científicas evidenciam tendências e influências nas mais diversas áreas do conhecimento, uma vez que são agentes de mudança na ciência e, consequentemente, na compreensão científica. Diante da necessidade de pesquisar como os campos dos conhecimentos sobre intenção empreendedora e sustentabilidade estão se estruturando no âmbito acadêmico, principalmente, considerando as análises de redes bibliométricas, assim como a análise de conteúdo (Knutas et al., 2015), questiona-se: como se estruturam as pesquisas científicas que abordam, conjuntamente, intenção empreendedora e sustentabilidade? Nesse sentido, delimitou-se como objetivo da pesquisa, portanto, investigar a produção científica internacional sobre intenção empreendedora e sustentabilidade. 

% Todavia, por meio de buscas sistêmicas em grandes portais de pesquisas científicas, como Spell, Scielo, Scopus e Web of Science, não foram constatados estudos sobre esses temas concomitantemente com caráter bibliométrico, adiante revistos, razão esta porque se entende que esta pesquisa pode suplantar uma lacuna da literatura científica e, a partir disso, contribuir diretamente para o desenvolvimento dos temas investigados.
% A inovação, a propagação da inovação e o surgimento de novos empreendimentos, em muitos países, são importantes sinais para o crescimento e recuperação de crises econômicas \cite{Silva2017EducacaoUFF}. A inovação é orientada de acordo com várias racionalidades, podendo ser observada por diversas óticas e utilizando diversos instrumentos para aprendizado \cite{Munoz2016LaProfesores}. Com efeito, o ambiente acadêmico se apresenta como uma unidade básica para o desenvolvimento de novos processos inovadores onde tais conteúdos devem ser amplamente explorados \cite{Costa2017InovacaoEnsino/aprendizagem}.
% 
% Desta forma, o desenvolvimento de uma visão empreendedora como ferramenta a inovação é essencial para formação de profissionais que tenham iniciativa, visão estratégica e capacidade de liderança, perfil este requerido pelas grandes empresas e startups \citeonline{Garcia2000FormacaoEmpreendedorismo}. Além das características já citadas, é imprescindível para o empreendedor desenvolver sua criatividade, pois assim, torna-se mais fácil encontrar ideias originais e com valor \cite{Macedo2019CapitalIdeologia}.
% 

 
 


TEXTO SOBRE JOGO 

Mostramos como as experiências de jogo dos consumidores através do flow aumentam o engajamento com um aplicativo gamificado, e como isso resulta em uma melhoria nas percepções de valor na realização de um comportamento de energia sustentável. Além disso, evidenciamos como o valor no comportamento criado por um aplicativo gamificado não apenas influencia as intenções comportamentais de realizar um comportamento de energia sustentável, mas também as intenções de continuar usando o aplicativo gamificado. Estes resultados fornecem importantes insights teóricos e práticos sobre o potencial da gamificação para ser usada em marketing sustentável e como as experiências lúdicas (flow e engajamento do cliente) podem ser transferidas para percepções positivas de comportamentos sustentáveis (valor no comportamento).


TEXTO SOBRE SUSTENTABILIDADE

Conforme esta categoria de política ambiental, estão os Objetivos de Desenvolvimento Sustentável (ODS) estabelecidos pela FAO (FAO, 2014). As estratégias para alcançar os ODS se baseiam em abordagens integradas e multissetoriais, incluindo a participação ampla de diferentes atores, sejam eles públicos ou privados, visando promover sinergia entre setores como agricultura, alimentação, saúde, nutrição, educação, desenvolvimento social, economia, entre outros. Dessa forma, é necessário que o enfrentamento da degradação florestal e da exploração não sustentável das florestas seja discutido em conjunto com a agricultura, seus sistemas, a silvicultura e as atividades comerciais, para que estes possam ser tornados renováveis (Silva, 2018).


Considerando tais fatores, este trabalho relaciona-se estritamente com o Desenvolvimento Rural de forma sustentável e Circular. Em tempos passados, o conceito de Desenvolvimento Rural Sustentável era denominado por “Progresso Rural”, pois, havia um entendido genérico como sentido parcial e prático de “Melhoramento do ambiente” \cite{almeida_da_1995}. Entretanto, torna-se imprescindível destacar que, o desenvolvimento sustentável no meio rural não pode ter suas bases de compreensão apenas no progresso econômico, local ou regional.



TEXTO SOBRE EMPREENDEDORISMO 

% A inovação, a propagação da inovação e o surgimento de novos empreendimentos, em muitos países, são importantes sinais para o crescimento e recuperação de crises econômicas \cite{Silva2017EducacaoUFF}. A inovação é orientada de acordo com várias racionalidades, podendo ser observada por diversas óticas e utilizando diversos instrumentos para aprendizado \cite{Munoz2016LaProfesores}. Com efeito, o ambiente acadêmico se apresenta como uma unidade básica para o desenvolvimento de novos processos inovadores onde tais conteúdos devem ser amplamente explorados \cite{Costa2017InovacaoEnsino/aprendizagem}.
% 

% Desta forma, o desenvolvimento de uma visão empreendedora como ferramenta a inovação é essencial para formação de profissionais que tenham iniciativa, visão estratégica e capacidade de liderança, perfil este requerido pelas grandes empresas e startups \citeonline{Garcia2000FormacaoEmpreendedorismo}. Além das características já citadas, é imprescindível para o empreendedor desenvolver sua criatividade, pois assim, torna-se mais fácil encontrar ideias originais e com valor \cite{Macedo2019CapitalIdeologia}.
% 
% 
% Para a inovação é necessário autoeficácia, pensar e agir diferente, encontrar soluções alternativas para os problemas e buscar ideias que tragam melhorias. Quando se fomenta a criatividade no ambiente educacional, a conquista da autonomia é consequência, assim como também a adoção de uma postura empreendedora \cite{Gonzalez2009PredictorsClasses}. Entretanto, na maioria das vezes, temos observado que essa importante característica não está sendo adequadamente desenvolvida no meio acadêmico. 



%%%%%%%%%%Antigos%%%%%

TEXTO SOBRE SUSTENTABILIDADE E MANEJO FLORESTAL


Inicialmente, as ações ambientais eram destinadas às empresas com maior suporte financeiro, dessa forma os setores responsáveis se afastavam das demandas que permeavam a estrutura social do país, repassando as responsabilidades para os demais. Atualmente, as ações ambientais são vistas como importantes fatores de contribuição para o desenvolvimento sustentado de um país. Se de um lado, a Revolução Industrial apresentou um novo cenário para as atividades econômicas e humanas, os novos processos, tecnologias e produtos transformaram as concepções sobre os recursos naturais, o ser humano e o Planeta. 


Para o desenvolvimento sustentável, a ciência e a tecnologia correspondem a um sistema de articulação entre uma racionalidade ambiental do processo de desenvolvimento e os processos concretos que definem as possibilidades de estratégias de manejo integrado do meio ambiente. Essa interação requer que o sistema de ciência e tecnologia esteja sustentado por paradigmas que incorporem o potencial ecológico, as condições ambientais e os valores culturais na organização dos processos produtivos \cite{Furlan2018}.

O desenvolvimento sustentável e a responsabilidade social assumem um papel cada vez mais importante nas estratégias das empresas. O agravamento dos problemas ambientais está relacionado a maneira como o conhecimento técnico-científico tem sido aplicado no processo produtivo. Os danos ao meio ambiente não são fatos inesperados e imprevisíveis, apenas demonstram a falta de capacidade do conhecimento de controlar os efeitos gerados pelo desenvolvimento industrial \cite{Maranhao2016}. 


Esse processo de planejamento sustentável inclui monitoramento e avaliação das ações realizadas por manejo florestal sustentável. 
A compreensão da sustentabilidade ambiental seja ela para fins acadêmicos, ecológicos ou econômicos requer também a compreensão do manejo florestal sustentável. Podemos definir o manejo florestal, em seu sentido mais amplo, como o conjunto de medidas tomadas em relação à floresta, principalmente de caráter silvocultural, visando otimizar a produção de determinados bens e/ou serviços de forma sustentável ao longo do tempo \cite{Rosot2010ManejoAraucaria}.


\citeonline{Louman2004PlanificacionTropicales} estenderam esse conceito para o “manejo florestal diversificado” que, à semelhança do ordenamento territorial, busca a ótima combinação dos recursos florestais, considerando o ponto de vista de seus proprietários e/ou usuários. A extensão de conceito enfatiza a importância da silvicultura de base comunitária para garantir a segurança alimentar e consonância com a proteção ambiental.

Para \citeonline{Teixeira2019a} as políticas públicas — tidas como diretrizes e princípios norteadores de ação do poder público — tem caráter estratégico para a condução de uma nação. Quando associadas à CT\&I, essas políticas são fatores determinantes para o crescimento econômico de um país, especialmente no que diz respeito à capacidade de produção de bens, serviços e a influência no padrão de vida de sua população \cite{Corazza2004}. Programas, políticas públicas e estratégias econômicas que capacitem os provedores de serviços técnicos, bem como produtores e comunidades locais, exercem impactos significativos na conservação das unidades florestais assim como o crescimento do mercado de produtos florestais sustentáveis cultivados em florestas \cite{Trozzo2021ForestIn}. O investimento em empreendedorismo sustentável e bioeconomia é um exemplo de  lastro político e gerencial em busca do alcance dos objetivos do desenvolvimento sustentável.


A empresa inovadora assim como o ser humano inovador está sempre em busca de soluções, e consegue enxergar oportunidades nos problemas e riscos previstos, tendo iniciativa e sendo proativa e visionária \cite{Milian2020}. Dessa forma, podendo contribuir para introduzir inovações na empresa em que trabalha, ajudar a solucionar problemas da sua cidade e no campo, entre outras possibilidades \cite{Alencar2019,Loiola2016}. Segundo \citeonline{DeMori1998} fatores condicionantes do sucesso no empreendedorismo autores como fundamentaram seus estudos em três grupos: características individuais do fundador, características estruturais e estratégicas do novo negócio, e condições determinadas pelo meio ambiente da empresa. A busca por inovação de processos de produção e no gerenciamento de vendas possibilita o alcance dos níveis mais elevados de competitividade. 

Ainda tratando sobre empreendedorismo sustentável, a literatura especializada reporta que o desenvolvimento econômico de um país depende de uma integração qualificada de empreendedores, sustentabilidade ambiental e da economia mundial \cite{Swinburn2006}. Para os empreendedores serem bem-sucedidos em mercados internacionais são necessárias vantagens competitivas baseadas em inovação, complexidade e/ou sofisticação dos produtos e um complexo sistema de compensação dos impactos ambientais que causam, isto é, devem ser inovadoras ou intensivas em conhecimento da economia circular \cite{Lucas2019}. O relatório Brundtland define o desenvolvimento sustentável como um progresso que atende às necessidades do presente sem comprometer a capacidade das gerações futuras de atender às suas próprias necessidades \cite{Brundtland1987RelatorioBrundtland}.

Porém, com as ações do homem agredindo a natureza em quase todos os processos naturais, a união entre o empreendedorismo com o impulso empreendedorismo e a sustentabilidade ambiental parece ser um “casamento entre valores opostos”. No entanto, é precisamente essa combinação que os formuladores de políticas contemporâneos promovem como a novidade do empreendedorismo socioambiental \cite{Mutyasira2020,Zahra2009AChallenges}. 

A perspectiva de mudanças climáticas emergentes requer que a perspectiva da ecologia do manejo florestal mude do âmbito local e regional para o global \cite{su12145667}. A implementação efetiva de práticas de manejo florestal sustentável depende da realização das operações florestais de maneira sustentável \cite{Marchi2018SustainableClimate}. Investimento suficiente em florestas e seus povos tradicionais devem ser usados para garantir a manutenção contínua das florestas.

Diferentes formas de explorações econômicas sustentáveis rurais e florestais têm sido pesquisadas, produzidas e utilizadas nos últimos anos, usufruindo das paisagens naturais sem agressão ao meio ambiente, como o já citado  Manejo florestal comunitário e empreendimentos agrícolas familiares que exploram, ou têm o compromisso de explorar a natureza selvagem, e considerando o respeito ao meio ambiente \cite{BICHEL2021127714,doi:10.1080/00167428.2020.1798764}. Em outras situações como a atividade artesanal, preocupa-se com retira da natureza de seus insumos (matéria-prima) e recursos energéticos, preocupando-se com a manutenção e recuperação natural destes, buscando associar a produção sustentável e uma posição competitiva, promovendo uma governança mais eficiente dos recursos renováveis de base biológica, e social \cite{DAmato2020TowardsSMEs}. 

Desta forma tais empresas associando-se perfeitamente a bioeconomia circular, estando alinhadas perfeitamente a agenda 2030 dos Objetivos de Desenvolvimento Sustentável da ONU \cite{Mundo2016}.


TEXTO SOBRE MARCA E PROPRIEDADE INTELECTUAL 

A Propriedade Intelectual possui um conceito amplo que introduz temáticas relacionadas ao processo de proteção às criações da mente, envolvendo aspectos relacionados à Propriedade Industrial, Direito Autoral e Direitos Conexos. Observa-se que há uma preocupação em relação à proteção a referida propriedade, como constatado nos principais tratados e acordos a seguir: Convenção da União de Paris (CUP), Tratado de Madri, Convenção de Berna, Convenção de Novas Variedades de Plantas, Acordo Geral de Tarifa e Comércio, Acordo TRIPS (\textit{Trade Related Aspects of Intellectual Property Rights}) e a Rodada Uruguai \cite{FERRAZ2008}. No Brasil o Instituto Nacional da Propriedade Industrial (INPI), criado em 1970, é uma autarquia federal vinculada ao Ministério da Indústria, Comércio Exterior e Serviços (MDIC), responsável pela execução das normas que regulamentam a propriedade intelectual para indústria no Brasil em âmbito social, econômico, jurídico e técnico, cabendo ainda o pronunciamento quanto ao interesse de assinatura, ratificação e denúncia de convenções, tratado, convênios, assim como acordos sobre a referida propriedade \cite{INPI2020}, destacando-se a “marca” como relevante fator para distinguir um produto ou serviço entre organizações, tendo o \textit{brand equity} um definidor de valor comercial a este ativo empresarial intangível.

Em busca de uma definição que abranja todas as possíveis facetas que descrevem uma marca (marketing, jurídica e comportamental), Feldwick (1996) sistematiza várias interpretações de brand equity a partir da tipologia tripartida, de modo a esclarecer discussões em torno desses conceitos. Os conceitos, contidos em um brand equity, se construídos de maneira adequada, promovem a imaginação dos consumidores e convencem sobre a marca, formando pistas sensoriais para ajudá-los a imaginar que estão comprando produtos confiáveis, nesse caso produtos ambientalmente sustentáveis, formando imagens mentais baseadas em informações pré-existentes e novas informações contidas na identidade visual das marcas (Cowan & Dai, 2014; Xie et al., 2020). 

Uma marca registrada é um identificador único para um produto, empresa ou serviço. É legalmente protegido por meio de registro de solicitação nos órgãos públicos de cada nação. Esse processo torna as marcas registradas importantes ativos econômicos para os registrantes. As marcas enquadram-se na subárea “Propriedade Industrial”, definida como um conjunto de elementos tangíveis e simbólicos que influenciam e determinam as preferências do consumidor. Segundo \cite{Mourad2014}, essa definição decorre do fato de as marcas conterem valor econômico associado. Fora do período de validade geral de dez anos, as marcas registradas podem ser prorrogadas, podendo ser realizada a renovação pelo titular da marca a seu pedido, com resultados iguais para cada extensão.

Para a inovação é necessário autoeficácia, pensar e agir diferente, encontrar soluções alternativas para os problemas e buscar ideias que tragam melhorias. Quando se fomenta a criatividade no ambiente educacional, a conquista da autonomia é consequência, assim como também a adoção de uma postura empreendedora \cite{Gonzalez2009PredictorsClasses}. Entretanto, na maioria das vezes, observamos que essa importante característica não está sendo adequadamente desenvolvida no meio acadêmico. 



%%%MANEJO FLORESTAL 

O manejo florestal sustentável inclui monitoramento e avaliação das ações e é essencial para a compreensão da sustentabilidade ambiental. O manejo florestal é o conjunto de medidas para otimizar a produção de bens e serviços de forma sustentável (Rosot, 2010). No Brasil, estratégias como gerenciamento adequado de resíduos, uso de fontes energéticas renováveis e produção extrativista estão sendo desenvolvidas para minimizar impactos ambientais na floresta. Essas técnicas ajudam a melhorar os índices sociais, econômicos e ambientais relacionados à sustentabilidade. Diferentes formas de exploração econômica sustentável rural e florestal têm sido pesquisadas e usadas recentemente, aproveitando paisagens naturais sem prejudicar o meio ambiente (Bichel & Telles, 2021; Karnani, 2011). 